% =====================================================================
\section{Evaluation, limitations and conclusions}\label{sec:evaluation}
% =====================================================================

\subsection{Later outcomes}

The plan and the evaluation measures were frozen on 17 November 2026, before any 2025/26 record was opened, but by people who already knew parts of what happened in 2025, so the evaluation describes outcomes and is not a prospective test (Appendix~\ref{app:history}). Table~\ref{tab:evaluation} covers the six purchases, the eighteen direct alternatives, United's five actual permanent senior arrivals (Table~\ref{tab:arrivals}) and the three players the plan sells, using the share of the main club's league minutes and the change in market value to 12 June 2026. All six purchases played between 0.50 and 0.78 of their main club's league minutes. Five later played for a different club than at the cutoff, Mbeumo among them at United, but these moves do not show that United could have bought them at the assumed terms: Hincapié joined Arsenal on a loan made permanent only in 2026 \citep{arsenal2026hincapie}, and David left Lille at the end of a contract \citep{bein2025david}. The three players the plan sells stayed at United and played 48\%, 76\% and 95\% of the league minutes, a risk indicator for those sales rather than proof that they were wrong. Nor does the comparison show that the model's purchases did better than United's actual arrivals, who have the higher mean share of league minutes (0.70 against 0.62) and the higher median change in value ($+$25.0\% against $-$13.3\%). The groups are small and unmatched, three members have no league minutes in the covered leagues, and the outcomes concern individual players in different teams, so no inferential test is used and no causal claim is made (Appendix~\ref{app:validation}). The outcomes are context for the proxy and the assumptions, not evidence of better recruitment.

\begin{table}[!htbp]
\centering
\caption{Outcomes in 2025/26 by group. ``Valid'' counts members with an observation: league minutes in the covered leagues, or a market value. $^{*}$Added after the evaluation as an extra description.}\label{tab:evaluation}
\tablefile{evaluation}
\end{table}

\subsection{Limitations}

\begin{itemize}
\item \emph{The score is a proxy.} It lacks defensive and goalkeeping statistics, weights market value although value also sets the fee, and can overstate the playing time of players who changed club during a season. An audit identifies 153 of the 830 reference players with league appearances for more than one club in a season, including the retained player Patrick Dorgu; none of the six purchases or three proposed sales appears in this group. The audit does not quantify the resulting measurement differences or establish whether recalculating the common percentiles would change the plan.
\item \emph{Finances are estimated.} Fees are multiples of value, candidate wages are assumed bands and United's wages come from a later table; the protected limit covers only candidate wages, with an assumed 10\% deviation.
\item \emph{Injury risk, development, sides of the pitch and other formations are not modelled.}
\item \emph{Deals are assumed to complete.} Willingness to sell and contract terms are not modelled, the deal labels are judgement and the shortlist evidence is incomplete.
\item \emph{The study is retrospective.} Some rules followed provisional plans, and the correction and the additional analyses came after the outcomes were examined. The correction changed no plan, and repeating the base case in the additional runs reproduced the original recommendation.
\end{itemize}

\subsection{Conclusion}

Under the stated assumptions, the model recommends buying Ritsu Doan, Bryan Mbeumo, Lucas Chevalier, Jonathan David, Éderson and Piero Hincapié and selling Harry Maguire, Casemiro and Luke Shaw. The plan spends \pounds 198.691m net, uses \pounds 147.94m of the \pounds 148m wage allowance and raises the quality proxy from 876.466 to 1,100.959 points. Under the assumed wage bands, wage allowance and arrival limit, the plan is limited by arrival places and wage room rather than by the transfer allowance: one more arrival is worth 20.5 points and \pounds 7m more wages 6.9, while \pounds 50m more transfer allowance is worth nothing, and the dual prices of the linear relaxation point the same way. This describes the planning scenario, not the constraints United actually faced. Doan, David and Chevalier are stable across the tested scenarios and are bought at every level of wage protection, Mbeumo is kept in most runs, and Hincapié and Éderson depend on the assumptions but are cheap to replace, even together. Protecting the plan fully against candidates earning 10\% more than estimated costs 7.5 points. The later outcomes show that the purchases received regular minutes, but not that they would have done better than United's actual arrivals, and they cannot confirm the plan. Future work could use dated wage and contract data, injury records for a non-neutral availability factor, and uncertainty sets for fees and values as well as wages.
